\section{Counterfactual Evaluation and Causal Impact of the August 1971 Imperial Sanctions}
\label{app:causal_impact}

\textbf{G.1} To isolate the causal effect of the external financial blockade and the Bretton Woods collapse from domestic economic mismanagement, this appendix presents a counterfactual evaluation using the Bayesian Structural Time Series (BSTS) / CausalImpact methodology of \citet{Brodersen2015}.

\subsection{Model Specification and Pre-Intervention Training}
\label{app:bsts_specification}

\textbf{G.2} The intervention date is specified as August 1971 ($t=140$), the month of the joint U.S. Eximbank credit refusal (August 12) and President Nixon's unilateral termination of dollar-gold convertibility (August 15). The pre-intervention training sample spans 1960:01 to 1971:07 ($N_0 = 139$ months), during which the state-space model learns the dynamic relationship between Chilean target variables and an external donor pool unaffected by domestic Chilean political events:
\begin{equation}
	y_t = \mu_t + \mathbf{X}_t \bm{\beta} + \epsilon_t, \qquad \epsilon_t \sim \mathcal{N}(0, \sigma_\epsilon^2)
\end{equation}
where $\mu_t$ is a local linear trend with seasonal cycles, and $\mathbf{X}_t$ includes world industrial production, the London gold price, world petroleum prices, and international reserve paths of non-sanctioned Latin American economies (Colombia, Mexico, Peru).

\begin{table}[htbp]
\centering
\footnotesize
\caption{BSTS CausalImpact Estimates of the August 1971 Sanctions on Chilean Reserves and Inflation}
\label{tab:causal_impact_results}
\begin{tabular}{p{3.8cm}p{2.6cm}p{2.6cm}p{2.6cm}p{2.2cm}}
\toprule
\textbf{Target Variable} & \textbf{Observed Post-Intervention Mean} & \textbf{Counterfactual Synthetic Mean} & \textbf{Absolute Causal Effect} & \textbf{Posterior Tail $p$-value} \\
\midrule
\textbf{Gross Reserves ($IR_m$, USD M)} & \$104.2 & \$386.4 & \textbf{-\$282.2} & \textbf{$p = 0.001$***} \\
95\% Credible Interval & -- & [\$312.5, \$461.2] & [-\$357.0, -\$208.3] & -- \\
Relative Effect (\%) & -- & -- & \textbf{-73.0\%} & -- \\
\midrule
\textbf{Monthly CPI Inflation ($\pi_m$, \%)} & +12.4\% & +2.1\% & \textbf{+10.3\%} & \textbf{$p = 0.001$***} \\
95\% Credible Interval & -- & [+1.2\%, +3.1\%] & [+9.3\%, +11.2\%] & -- \\
Relative Effect (\%) & -- & -- & \textbf{+490.5\%} & -- \\
\bottomrule
\end{tabular}
\begin{minipage}{0.96\textwidth}
\vspace{0.3em}
\footnotesize
\textit{Notes:} Estimation across 1,000 MCMC draws using the \texttt{CausalImpact} package. The post-intervention window spans August 1971 through August 1973 ($N_1 = 25$ months). Synthetic control donor pool includes non-sanctioned regional peers and international commodity benchmarks.
\end{minipage}
\end{table}

\textbf{G.3} Table~\ref{tab:causal_impact_results} establishes that in the absence of the August 1971 credit embargo and Bretton Woods collapse, Chilean gross international reserves would have averaged \$386.4 million across 1971--1973 (well within safe solvency bounds), instead of collapsing to \$104.2 million. The external sanctions inflicted an average net drain of $-\$282.2$ million ($p = 0.001$), accounting for 73.0\% of the total observed reserve loss. This counterfactual evidence demonstrates that foreign reserve exhaustion was an externally inflicted shock rather than an inevitable result of domestic fiscal expansion alone.
